\begin{center}
\SourceLeaf{026}\pdfbookmark[2]{8.}{article-008}8.
\end{center}
Eas iam inter anomalias atque radium vectorem relationes, quae imprimis attentione dignae sunt, colligamus, quarum deductio nemini in analysi trigonometrica vel mediocriter versato difficultates obiicere poterit. Pluribus harum formularum concinnitas maior conciliatur, introducto pro $e$ angulo cuius sinus est $=e$. Quo per $\varphi$ designato, habemus $\BareRoot{(1-ee)}=\cos\varphi$, $\BareRoot{(1+e)}=\cos(45^0-\frac{1}{2}\varphi)\BareRoot{2}$, $\BareRoot{(1-e)}=\cos(45^0+\frac{1}{2}\varphi)\BareRoot{2}$, $\BareRoot{\frac{1-e}{1+e}}=\operatorname{tang}(45^0-\frac{1}{2}\varphi)$, $\BareRoot{(1+e)}+\BareRoot{(1-e)}=2\cos\frac{1}{2}\varphi$, $\BareRoot{(1+e)}-\BareRoot{(1-e)}=2\sin\frac{1}{2}\varphi$. Ecce iam relationes praecipuas inter $a$, $p$, $r$, $e$, $\varphi$, $v$, $E$, $M$.\HistoricalNoteEight{Handschriftliche Bemerkung:]
\[
\begin{gathered}
1=\cos v\cos E+\frac{\sin v\sin E}{\cos\varphi},\\
\cos\varphi=\frac{\sin v\sin E}{1-\cos v\cos E},\\
\sin\varphi=\frac{\cos E-\cos v}{1-\cos v\cos E}\cdot
\end{gathered}
\]
}
\[
\mathrm{I.}\quad p=a\cos\varphi^2
\]
\[
\mathrm{II.}\quad r=\frac{p}{1+e\cos v}
\]
\[
\mathrm{III.}\quad r=a(1-e\cos E)
\]
\[
\mathrm{IV.}\quad \cos E=\frac{\cos v+e}{1+e\cos v},\quad\mathrm{sive}\quad\cos v=\frac{\cos E-e}{1-e\cos E}
\]
\[
\begin{aligned}
\mathrm{V.}\quad\sin\frac{1}{2}E&=\BareRoot{\frac{1}{2}(1-\cos E)}\\
&=\sin\frac{1}{2}v\BareRoot{\frac{1-e}{1+e\cos v}}=\sin\frac{1}{2}v\BareRoot{\frac{r(1-e)}{p}}\\
&=\sin\frac{1}{2}v\BareRoot{\frac{r}{a(1+e)}}
\end{aligned}
\]
\[
\begin{aligned}
\mathrm{VI.}\quad\cos\frac{1}{2}E&=\BareRoot{\frac{1}{2}(1+\cos E)}\\
&=\cos\frac{1}{2}v\BareRoot{\frac{1+e}{1+e\cos v}}=\cos\frac{1}{2}v\BareRoot{\frac{r(1+e)}{p}}\\
&=\cos\frac{1}{2}v\BareRoot{\frac{r}{a(1-e)}}
\end{aligned}
\]
\[
\mathrm{VII.}\quad\operatorname{tang}\frac{1}{2}E=\operatorname{tang}\frac{1}{2}v\operatorname{tang}(45^0-\frac{1}{2}\varphi)
\]
\[
\mathrm{VIII.}\quad\sin E=\frac{r\sin v\cos\varphi}{p}=\frac{r\sin v}{a\cos\varphi}
\]
\[
\begin{aligned}
\mathrm{IX.}\quad r\cos v&=a(\cos E-e)\\
&=2a\cos(\frac{1}{2}E+\frac{1}{2}\varphi+45^0)\cos(\frac{1}{2}E-\frac{1}{2}\varphi-45^0)
\end{aligned}
\]
\[
\begin{aligned}
\mathrm{X.}\quad\sin\frac{1}{2}(v-E)&=\sin\frac{1}{2}\varphi\sin v\BareRoot{\frac{r}{p}}\\
&=\sin\frac{1}{2}\varphi\sin E\BareRoot{\frac{a}{r}}
\end{aligned}
\]
\SourceLeaf{027}
\[
\begin{aligned}
\mathrm{XI.}\quad\sin\frac{1}{2}(v+E)&=\cos\frac{1}{2}\varphi\sin v\BareRoot{\frac{r}{p}}\\
&=\cos\frac{1}{2}\varphi\sin E\BareRoot{\frac{a}{r}}
\end{aligned}
\]
\[
\mathrm{XII.}\quad M=E-e\sin E.
\]
